\documentclass[tikz,border=3mm,multi=tikzpicture]{standalone}
\usepackage{eleve-matematica-ef1}

\begin{document}

% FIGURA 1 — fig-01-translacao-na-malha
\begin{tikzpicture}[matematica figura, x=.85cm, y=.85cm]
  \path[use as bounding box] (-0.30,-0.25) rectangle (11.10,8.10);
  \draw[step=1,matematica divisao] (0,0) grid (10,7);
  \fill[matemazulclaro] (1,1)--(3,1)--(2,3)--cycle;
  \draw[matematica contorno,line width=2pt] (1,1)--(3,1)--(2,3)--cycle;
  \fill[matemalaranjaclaro] (5,3)--(7,3)--(6,5)--cycle;
  \draw[matematica destaque,line width=2pt] (5,3)--(7,3)--(6,5)--cycle;
  \draw[-{Stealth[length=3mm]},matematica destaque] (2.25,2.10)--(5.60,4.10);
  \node[matematica resultado,fill=white,inner sep=3pt] at (4.25,5.45) {$4$ à direita};
  \node[matematica resultado,fill=white,inner sep=3pt] at (4.25,4.75) {$2$ para cima};
\end{tikzpicture}

% FIGURA 2 — fig-02-rotacao-na-malha
\begin{tikzpicture}[matematica figura, x=1cm, y=1cm]
  \path[use as bounding box] (-0.30,-0.20) rectangle (9.30,8.40);
  \draw[step=1,matematica divisao] (0.50,0.50) grid (8.50,7.50);
  \fill[matemacinza] (4.50,4.00) circle (.13);
  \draw[-{Stealth[length=3mm]},matematica contorno,line width=2.8pt] (4.50,4.00)--(7.20,4.00);
  \draw[-{Stealth[length=3mm]},matematica destaque,line width=2.8pt] (4.50,4.00)--(4.50,6.70);
  \draw[-{Stealth[length=3mm]},draw=matemaverde,line width=2.8pt] (4.50,4.00)--(1.80,4.00);
  \draw[-{Stealth[length=2.5mm]},matematica destaque] (6.20,4.45) arc[start angle=15,end angle=78,radius=1.78];
  \draw[-{Stealth[length=2.5mm]},draw=matemaverde,line width=1.8pt] (5.90,3.45) arc[start angle=-20,end angle=-160,radius=1.48];
  \node[matematica resultado] at (6.55,6.20) {$90^\circ$};
  \node[font=\Large\bfseries,text=matemaverde] at (2.20,2.40) {$180^\circ$};
  \node[matematica rotulo] at (4.50,3.35) {centro};
\end{tikzpicture}

% FIGURA 3 — fig-03-reflexao-na-malha
\begin{tikzpicture}[matematica figura, x=.90cm, y=.90cm]
  \path[use as bounding box] (-0.30,-0.25) rectangle (10.80,7.30);
  \draw[step=1,matematica divisao] (0,0) grid (10,6);
  \draw[matematica destaque,dashed,line width=2pt] (5,0)--(5,6);
  \fill[matemazulclaro] (1,1)--(4,1)--(4,4)--(3,4)--cycle;
  \draw[matematica contorno,line width=2pt] (1,1)--(4,1)--(4,4)--(3,4)--cycle;
  \fill[matemalaranjaclaro] (9,1)--(6,1)--(6,4)--(7,4)--cycle;
  \draw[matematica destaque,line width=2pt] (9,1)--(6,1)--(6,4)--(7,4)--cycle;
  \draw[matematica auxiliar] (4,4)--(6,4);
  \node[matematica valor,fill=white,inner sep=2pt] at (4.50,4.35) {$1$};
  \node[matematica valor,fill=white,inner sep=2pt] at (5.50,4.35) {$1$};
  \node[matematica resultado,rotate=90,fill=white,inner sep=3pt] at (5,3) {eixo};
\end{tikzpicture}

% FIGURA 4 — fig-04-ampliacao-na-malha
\begin{tikzpicture}[matematica figura, x=.72cm, y=.72cm]
  \path[use as bounding box] (-0.30,-0.25) rectangle (15.30,9.00);
  \draw[step=1,matematica divisao] (0,0) grid (5,7);
  \fill[matemazulclaro] (1,1) rectangle (3,4);
  \draw[matematica contorno,line width=2pt] (1,1) rectangle (3,4);
  \node[matematica valor] at (2,7.75) {$2\times3$};
  \draw[-{Stealth[length=3mm]},matematica destaque] (5.45,3.50)--(7.30,3.50);
  \node[matematica resultado] at (6.35,4.10) {dobrou};
  \draw[step=1,matematica divisao] (7.80,0) grid (14.80,7);
  \fill[matemalaranjaclaro] (9,0.50) rectangle (13,6.50);
  \draw[matematica destaque,line width=2pt] (9,0.50) rectangle (13,6.50);
  \node[matematica valor] at (11,7.75) {$4\times6$};
\end{tikzpicture}

% FIGURA 5 — fig-05-reducao-na-malha
\begin{tikzpicture}[matematica figura, x=.72cm, y=.72cm]
  \path[use as bounding box] (-0.30,-0.25) rectangle (15.30,8.70);
  \draw[step=1,matematica divisao] (0,0) grid (7,6);
  \fill[matemazulclaro] (0.50,1) rectangle (6.50,5);
  \draw[matematica contorno,line width=2pt] (0.50,1) rectangle (6.50,5);
  \node[matematica valor] at (3.50,7.00) {$6\times4$};
  \draw[-{Stealth[length=3mm]},matematica destaque] (7.50,3.00)--(9.30,3.00);
  \node[matematica resultado] at (8.40,3.65) {metade};
  \draw[step=1,matematica divisao] (9.80,0) grid (14.80,6);
  \fill[matemalaranjaclaro] (10.80,2) rectangle (13.80,4);
  \draw[matematica destaque,line width=2pt] (10.80,2) rectangle (13.80,4);
  \node[matematica valor] at (12.30,7.00) {$3\times2$};
\end{tikzpicture}

\end{document}
